% !TeX encoding = UTF-8
\documentclass[variant=docente,cover=institutional,mode=screen]{ua-engineering-v6-article}
% !TeX encoding = UTF-8
\UASetUniversity{Universidad Austral}
\UASetFaculty{Facultad de Ingeniería}
\UASetDepartment{Departamento de Computación}
\UASetCampus{Pilar}
\UASetCareer{Ingeniería Informática}
\UASetCommission{Comisión A}
\UASetProfessor{Equipo docente}
\UASetSemester{Segundo cuatrimestre 2026}
\UASetCourse{Sistemas Distribuidos}
\UASetDate{2026}

\UASetSemester{Segundo cuatrimestre 2026}
\UASetDate{2026}
\UASetClassNumber{07}
\UASetClassTitle{Logs distribuidos y procesamiento de eventos}
\title{Clase 07 --- Guía docente: de un hecho durable a un efecto recuperable}
\author{\UAProfessor}
\date{\UADate}
\begin{document}
\maketitle
\tableofcontents
\clearpage

\section{Parte A: guion operativo para 180 minutos}

\subsection{Objetivo pedagógico profundo}

La Clase 06 mostró que parte del grupo aún no posee un modelo firme de memoria, persistencia, archivos y progreso recuperable. Por eso esta clase no debe comenzar con Kafka, offsets o exactly-once. Debe construir primero un modelo mínimo:

\begin{center}
\textbf{registro} $\rightarrow$ \textbf{secuencia append-only} $\rightarrow$ \textbf{lector} $\rightarrow$ \textbf{progreso durable}.
\end{center}

Sobre esa base se incorporan particiones, keys, grupos de consumidores, fallas e idempotencia.

\begin{uakeybox}
La clase habrá funcionado si el estudiante puede mover un crash en un timeline y deducir pérdida o duplicación sin recitar de memoria una etiqueta de entrega.
\end{uakeybox}

La progresión será:
\[
\text{situación observable}
\rightarrow
\text{predicción}
\rightarrow
\text{objetos concretos}
\rightarrow
\text{intuición}
\rightarrow
\text{definición}
\rightarrow
\text{falla}
\rightarrow
\text{evidencia}
\rightarrow
\text{decisión}.
\]

\subsection{Resultados de aprendizaje observables}

Al finalizar, el estudiante debe poder:
\begin{itemize}
\item diferenciar memoria de trabajo, historia persistente y progreso durable;
\item explicar registro, append, log, partición, key, offset y committed offset;
\item elegir una key a partir del orden exigido por el dominio;
\item construir timelines de at-most-once y at-least-once;
\item justificar idempotencia e inbox para un efecto local;
\item descubrir las dos ventanas de dual write y explicar outbox;
\item declarar la frontera de una garantía exactly-once;
\item diagnosticar lag por partición;
\item ubicar late data según event time y una política explícita.
\end{itemize}

No es objetivo nuclear memorizar configuraciones de Kafka, ISR, transacciones internas o restauración avanzada de state stores. Esos temas son extensión o respaldo.

\subsection{Arquitectura didáctica v9 reestructurada para 180 minutos}

\begin{center}
\begin{tabular}{p{0.12\textwidth}p{0.22\textwidth}p{0.56\textwidth}}
\toprule
Tiempo & Bloque & Propósito \\
\midrule
0--12 & Caso real y predicción & CampusShop; distinguir operación, evento y efecto sin usar jerga. \\
12--30 & Base mínima & Memoria, persistencia, registro, append, lector y progreso durable. \\
30--48 & De RPC a historia compartida & Comparar llamadas sincrónicas con publicar un hecho durable. \\
48--68 & Log, partición y key & Tarjetas físicas; orden local y hot partition. \\
68--82 & Actividad 1 & Distribuir eventos y defender una key. \\
82--92 & Pausa y recuperación & Repasar vocabulario mediante preguntas breves. \\
92--112 & Consumer, offset y grupo & Position, committed offset y reparto de particiones. \\
112--138 & Timelines de falla & At-most-once, at-least-once, event ID e inbox. \\
138--156 & Dual write y outbox & Descubrir las dos ventanas antes de nombrar el patrón. \\
156--169 & Tiempo en streams & Ventana, event time y late data; watermark como extensión. \\
169--178 & Actividad 2 & CampusShop o Asteria: diseño y defensa. \\
178--180 & Exit ticket & Cuatro respuestas para abrir Clase 08. \\
\bottomrule
\end{tabular}
\end{center}

No debe haber más de doce o quince minutos continuos de exposición. Las diapositivas EXTENSIÓN y RESPALDO no forman parte de la ruta obligatoria.

\subsection{Preparación previa del docente}

\subsubsection*{Cuarenta y ocho horas antes}
\begin{itemize}
\item Imprimir tarjetas de eventos, tres carriles, dos consumidores, una ficha de crash y dos marcadores de progreso.
\item Ejecutar las demos de offsets, hot partition y outbox/inbox.
\item Ensayar el caso CampusShop en menos de noventa segundos.
\item Revisar operation ID, event ID y partición+offset.
\item Elegir un solo caso adicional: ranking de Asteria para event time.
\end{itemize}

\subsubsection*{Treinta minutos antes}
\begin{itemize}
\item Dibujar tres particiones P0--P2.
\item Escribir: \textbf{leer $\neq$ procesar $\neq$ producir efecto $\neq$ confirmar progreso}.
\item Preparar el timeline Broker--Consumer--Efecto--Offset.
\item Dejar visible la pregunta central.
\end{itemize}

\subsubsection*{Cinco minutos antes}
\begin{itemize}
\item Confirmar el medio de predicción individual.
\item Entregar o dejar disponibles las fichas de aula.
\item Recordar que EXTENSIÓN y RESPALDO se usan solo si el ritmo lo permite.
\end{itemize}

\subsection{Plan de pizarrón}

Divida el pizarrón en tres zonas permanentes:
\begin{enumerate}
\item \textbf{Historia:} particiones y registros.
\item \textbf{Progreso:} position y committed offset.
\item \textbf{Efecto:} estado de negocio, event ID y crash.
\end{enumerate}

Mantenga los timelines de pérdida y duplicación visibles al mismo tiempo. La comparación simultánea reduce carga cognitiva.

\subsection{Guion de apertura}

Diga aproximadamente:
\begin{quote}
``La clase anterior mostró cómo un nodo conserva una historia. Hoy esa historia se convierte en un punto de encuentro entre aplicaciones. CampusShop confirma una compra; inventario, email y analítica deben reaccionar, pero no queremos que la caída de email detenga checkout. Vamos a seguir un hecho y mover un crash por todo el recorrido.''
\end{quote}

Dé 45 segundos individuales y pregunte:
\begin{itemize}
\item ¿qué consumidor no puede duplicar su efecto?;
\item ¿cuál puede atrasarse?;
\item ¿qué debe recordar un consumidor después de caer?;
\item ¿qué dato identifica la misma operación lógica?
\end{itemize}

Registre hipótesis sin corregirlas todas. Se revisarán al final.

\subsection{Base mínima antes del vocabulario distribuido}

No suponga conocimientos sólidos de sistemas de archivos. Utilice solamente:
\begin{description}
\item[Memoria de trabajo.] Estado rápido que puede perderse al reiniciar.
\item[Registro.] Unidad que representa un hecho.
\item[Append.] Agregar al final de una secuencia.
\item[Log.] Historia retenida de registros.
\item[Lector.] Proceso que recorre esa historia.
\item[Posición actual.] Punto alcanzado por la ejecución viva.
\item[Progreso durable.] Marca desde la que puede continuar otra instancia.
\end{description}

Use la secuencia 416--420. Mueva el lector hasta 420, pero deje el progreso durable en 418. Pregunte qué reaparece tras un crash. Recién entonces introduzca \emph{position} y \emph{committed offset}.

\subsection{Secuencia de conducción: de la intuición al formalismo}

\subsubsection{De llamadas sincrónicas a un hecho durable}
Dibuje checkout llamando a inventario, email y analítica. Pregunte qué sucede si email tarda veinte segundos. Luego inserte un log. La conclusión debe ser:
\begin{itemize}
\item el log desacopla producción y procesamiento;
\item no elimina fallas ni define el efecto correcto;
\item desplaza complejidad hacia orden, progreso, idempotencia y operación.
\end{itemize}

\subsubsection{Log, partición y key}
Use tarjetas de dos pedidos. Cada carril es una partición. Pida una key antes de mostrar una fórmula. Si el grupo elige tipo de evento, muestre cómo un pedido puede perder orden. Luego introduzca:
\[
partition(key)=h(key)\bmod N.
\]
Pregunte por una guild o tenant caliente. La key decide simultáneamente orden y distribución.

\subsubsection{Productor y confirmación}
Explique el timeout ambiguo: el evento pudo no llegar o pudo quedar durable y perderse el ACK. Introduzca identidad estable y retry. Deje ISR y producer idempotence como profundización si el grupo todavía está consolidando la base.

\subsubsection{Consumer, position y committed offset}
Use el marcador físico. Position es volátil; committed offset es el contrato de recuperación del grupo. Haga un consumer group humano con tres particiones y dos estudiantes. Agregue un tercero y muestre que el paralelismo útil no supera la cantidad de particiones.

\subsubsection{Semánticas mediante timeline}
No defina at-most-once antes de la ejecución.

Timeline A: leer 418, confirmar 419, crash, efecto ausente. Nombre el daño y luego la semántica.

Timeline B: leer 418, aplicar efecto, crash antes de confirmar, releer y repetir. Nombre el daño y luego la semántica.

La deducción esperada es que no existe un orden mágico: para un recurso local hace falta event ID + efecto + inbox en la misma transacción.

\subsubsection{Exactly-once: declarar frontera}
No convierta exactly-once en el centro. Úselo para consolidar el método: exactamente una vez, ¿qué efecto, dentro de qué recursos y ante qué fallas? Un proveedor externo queda fuera si no participa del protocolo.

\subsubsection{Dual write y outbox}
Dibuje base y broker. Pruebe ambos órdenes con un crash intermedio. Espere que el grupo descubra que cambiar el orden no crea atomicidad. Introduzca outbox como negocio + evento pendiente en una transacción local. Haga perder el ACK del relay para mostrar por qué todavía se requiere inbox/idempotencia.

\subsubsection{Event time y late data}
Use un evento ocurrido a las 10:04:50 y recibido a las 10:07:10. Event time resuelve la pertenencia; la política de late data resuelve qué hacer con la llegada posterior. Si falta tiempo, no formalice watermarks: deje event time, processing time y tres políticas de late data.

\subsection{Demostraciones guiadas}

\begin{uainfo}
\textbf{Demo 1 --- secuencia y progreso.} Tarjetas 416--420, un lector y dos marcadores. Haga caer el lector y determine qué reaparece.
\end{uainfo}
\begin{uainfo}
\textbf{Demo 2 --- hot partition.} Reparta eventos con una key sesgada. Muestre que consumidores ociosos no dividen el carril caliente.
\end{uainfo}
\begin{uainfo}
\textbf{Demo 3 --- crash móvil.} Cuatro estudiantes representan broker, consumer, efecto y offset. Mueva la tarjeta CRASH.
\end{uainfo}
\begin{uainfo}
\textbf{Demo 4 --- outbox e inbox.} Pierda el ACK del relay, publique dos copias con el mismo event ID y produzca un solo efecto lógico.
\end{uainfo}

\subsection{Preguntas socráticas y respuestas esperadas}

\begin{tabular}{p{0.45\textwidth}p{0.45\textwidth}}
\toprule
Pregunta & Núcleo esperado \\
\midrule
¿Qué debe persistir un lector? & progreso suficiente para reiniciar \\
¿Event ID y offset son iguales? & identidad del hecho frente a posición \\
¿Dónde existe orden? & dentro de una partición \\
¿Por qué una key puede ser mala? & rompe orden o concentra carga \\
¿Qué pasa si efecto terminó y offset no? & replay y posible duplicación \\
¿Qué resuelve inbox? & efecto y event ID atómicos localmente \\
¿Qué resuelve outbox? & negocio e intención de publicar atómicos \\
¿Qué deja fuera exactly-once? & recursos no participantes \\
\bottomrule
\end{tabular}

\subsection{Errores previsibles y estrategias de corrección}

\begin{tabular}{p{0.36\textwidth}p{0.55\textwidth}}
\toprule
Error & Intervención \\
\midrule
``Kafka es una cola'' & Mostrar dos consumidores releyendo la misma historia. \\
``Offset es ID'' & Escribir offset 42 en P0 y P1. \\
``Más consumidores siempre ayuda'' & Comparar cinco consumidores con tres particiones. \\
``Lag cero significa correcto'' & Adelantar offset antes del efecto. \\
``At-least-once nunca pierde'' & Introducir un poison event o dependencia permanentemente fallida. \\
``Outbox elimina duplicados'' & Perder el ACK del relay. \\
``Exactly-once es global'' & Dibujar email fuera de la transacción. \\
``Event time es llegada'' & Usar un dispositivo offline. \\
\bottomrule
\end{tabular}

\subsection{Criterios de evaluación formativa}

Una respuesta lograda incluye:
\begin{enumerate}
\item hecho e identidad;
\item key y orden;
\item estado durable de progreso;
\item punto de crash;
\item efecto e idempotencia;
\item métrica o prueba.
\end{enumerate}

\begin{center}
\begin{tabular}{p{0.18\textwidth}p{0.70\textwidth}}
\toprule
Nivel & Evidencia \\
\midrule
Inicial & repite nombres sin dibujar ejecución \\
En desarrollo & identifica componentes, pero no frontera de falla \\
Logrado & timeline, garantía, límite y métrica \\
Avanzado & compara alternativas y propone prueba adversarial \\
\bottomrule
\end{tabular}
\end{center}

\subsection{Contingencias}

\subsubsection*{Si falta tiempo}
Preserve: base mínima, key/partición, committed offset, dos timelines, inbox y outbox. Reduzca ISR, producer idempotence, state stores, watermarks y DLQ.

\subsubsection*{Si la base del grupo sigue débil}
No agregue más definiciones. Repita el mismo evento con tarjetas y haga narrar qué existe en RAM, qué está persistido y desde dónde reinicia.

\subsubsection*{Si falla la demo}
La clase debe poder dictarse con tarjetas y pizarra; ningún concepto nuclear requiere un cluster Kafka.

\subsubsection*{Si hay poca participación}
Pida respuesta individual escrita antes de discutir en parejas. No comience con una pregunta abierta a toda el aula.

\subsection{Cómo favorecer una experiencia de clase excelente}

\begin{itemize}
\item Reconozca la brecha detectada en Clase 06 sin presentarla como déficit personal.
\item Haga visible la ruta mínima y no intente recorrer todo el respaldo.
\item Mantenga CampusShop como caso conductor; Asteria aparece solo como transferencia.
\item Defina pocos términos, úselos y recién después agregue otros.
\item Dibuje antes de nombrar la semántica.
\item Cierre cada bloque con ``qué resuelve, qué no resuelve y qué cuesta''.
\item Reserve ocho minutos para la actividad final y dos para el exit ticket.
\item Devuelva al inicio de Clase 08 una duda real del grupo.
\end{itemize}

\subsection{Conexión con el Trabajo Final Integrador}

Cada grupo debe completar para un flujo basado en eventos:
\begin{itemize}
\item hecho y event ID;
\item productor y owner;
\item key y orden;
\item consumidores y efectos;
\item política de progreso;
\item outbox/inbox;
\item SLO extremo a extremo;
\item crash o rebalance a ensayar.
\end{itemize}

No acepte ``usaremos Kafka porque escala''. Deben justificar desacoplamiento, replay u orden local.

\subsection{Cierre, exit ticket y puente}

Exit ticket:
\begin{enumerate}
\item Diferencia entre event ID y offset.
\item Timeline mínimo de duplicación.
\item Qué resuelve outbox y qué deja abierto.
\item Métrica para distinguir hot partition de consumidor lento.
\end{enumerate}

Cierre:
\begin{quote}
``Hoy vimos que la historia puede sobrevivir mientras distintos consumidores avanzan de manera diferente. En la próxima clase analizaremos qué ocurre cuando las réplicas mismas no pueden coordinarse.''
\end{quote}


\clearpage
\section{Parte B: dossier técnico de preparación}
\subsection{Dossier técnico de preparación del docente}
Esta sección concentra el conocimiento que conviene dominar para contestar preguntas durante la clase. No debe convertirse en exposición continua: la ruta esencial sigue siendo partición, offset, semánticas, frontera exactly-once y outbox.

\subsubsection{Modelo formal mínimo de log distribuido}

Un topic puede modelarse como un conjunto de particiones:
\[
T=\{P_0,P_1,\ldots,P_{m-1}\},\qquad
P_j=\langle e_{j,0},e_{j,1},\ldots\rangle.
\]
Cada partición ofrece un orden total local por offset. No existe necesariamente un orden total entre particiones. El offset es una posición dentro de una partición, no una identidad global del evento ni una prueba de que su efecto de negocio ocurrió.

La clave de partición cumple simultáneamente dos funciones:
\begin{itemize}
\item agrupa eventos que necesitan orden relativo;
\item distribuye carga y determina riesgo de hot partition.
\end{itemize}

Una elección correcta debe declarar la unidad de orden: pedido, cuenta, jugador, dispositivo o sesión. Elegir una clave de alta cardinalidad no basta si la distribución es sesgada.

\subsubsection{Publicación, replicación y ACK}

El producer envía un record al líder de la partición. La durabilidad observable depende de configuración y protocolo: cantidad de réplicas, conjunto de réplicas sincronizadas, política de ACK y tolerancia a fallas. El docente debe evitar la frase ``Kafka ya lo guardó'' sin declarar qué confirmación recibió el producer y frente a qué falla se considera durable.

La idempotencia del producer evita que retries del protocolo produzcan duplicados adicionales en el log dentro de la sesión y condiciones soportadas. No deduplica dos comandos de negocio distintos que usan IDs diferentes ni resuelve efectos externos del consumidor.

\paragraph{Pregunta previsible: ``¿orden del producer significa orden global?''}
No. El orden relevante se conserva por partición. Dos claves asignadas a particiones diferentes pueden observarse en órdenes distintos por consumidores diferentes.

\subsubsection{Consumer groups, position y committed offset}

En un consumer group, cada partición se asigna como máximo a un consumidor activo del grupo en un momento dado. Por eso el paralelismo útil de un grupo no excede la cantidad de particiones activas. Más consumidores que particiones quedan ociosos, aunque otros grupos pueden leer independientemente la misma historia.

Distinga:
\begin{itemize}
\item \textbf{position:} próximo offset que el consumidor local planea leer;
\item \textbf{committed offset:} punto durable desde el que el grupo retomará luego de reinicio o rebalance;
\item \textbf{efecto de negocio:} modificación externa producida por el record;
\item \textbf{event ID:} identidad semántica estable usada para deduplicar.
\end{itemize}

Un rebalance cambia propietarios de particiones. Si el consumidor no coordina su estado y offsets, puede repetir trabajo, perder cache caliente o aumentar latencia. El rebalance no es corrupción: es una transición operativa que exige diseño.

\subsubsection{Semánticas mediante timelines}

\paragraph{At-most-once.}
Se confirma el offset antes del efecto. Si hay crash entre ambos, el record no se reprocesará y el efecto puede faltar. ``Como máximo una vez'' describe intentos de entrega/procesamiento dentro de una frontera, no garantiza que el negocio suceda.

\paragraph{At-least-once.}
Se aplica el efecto y después se confirma el offset. Si hay crash en medio, el record se relee. Esta semántica permite redelivery; no garantiza que todo record produzca un efecto exitoso, porque puede existir poison pill, error permanente o indisponibilidad externa.

\paragraph{Exactly-once.}
Debe definirse como una propiedad de una frontera. En un pipeline que consume y vuelve a publicar dentro de Kafka, una transacción puede coordinar los offsets de entrada y los registros de salida. Un correo, un servicio de pago o una base externa no entra automáticamente en esa transacción. Para esos efectos suelen usarse idempotencia, outbox/inbox, claves únicas o una transacción local que incluya efecto y deduplicación.

\paragraph{Pregunta previsible: ``¿exactly-once significa que el código se ejecuta una sola vez?''}
No necesariamente. Puede haber reejecución física, pero el estado observable debe ser equivalente a una ejecución lógica. La distinción entre ejecución e impacto es central.

\subsubsection{Dual write, transactional outbox e inbox}

El dual write directo tiene dos ventanas:
\begin{enumerate}
\item la base confirma y la publicación falla: estado sin evento;
\item la publicación confirma y la base falla: evento que anuncia un estado inexistente.
\end{enumerate}

Transactional outbox persiste, en una misma transacción local, el cambio de negocio y una fila/evento pendiente de publicar. Un relay publica después y puede repetir; por eso el consumidor debe ser idempotente. Outbox resuelve atomicidad local entre estado e intención de publicar, no exactly-once global ni orden ilimitado entre agregados.

Inbox o tabla de deduplicación registra el event ID procesado junto con el efecto local. Una restricción única convierte la repetición en lectura del resultado anterior o no-op.

\subsubsection{Streaming stateful, changelog y recuperación}

Un operador stateless transforma cada evento sin memoria duradera. Un operador stateful mantiene agregados, ventanas, joins o sesiones. El estado debe tener:
\begin{itemize}
\item particionamiento compatible con las claves de entrada;
\item checkpoint o changelog recuperable;
\item offset consistente con el estado restaurado;
\item política de crecimiento, expiración y eventos tardíos.
\end{itemize}

Restaurar estado sin restaurar el offset coherente puede duplicar o saltar actualizaciones. Escalar un operador stateful implica redistribuir estado, no solo iniciar procesos.

\subsubsection{Event time, processing time y watermarks}

\begin{description}
\item[Event time.] Momento en que ocurrió el hecho en el dominio.
\item[Ingestion time.] Momento en que el sistema de streaming lo recibió.
\item[Processing time.] Momento en que un operador lo procesa.
\end{description}

Un watermark es una estimación de hasta qué event time se espera haber observado la mayoría de los eventos relevantes. No es una prueba absoluta de completitud. La política debe decidir qué hacer con late data: descartar, actualizar la ventana, emitir corrección o enviarlo a un flujo separado.

El costo de esperar más es mayor latencia y estado; cerrar temprano reduce latencia pero aumenta correcciones o pérdida semántica.

\subsubsection{Operación y métricas}

Las métricas mínimas incluyen:
\begin{itemize}
\item records/s y bytes/s por partición;
\item producer error/retry rate y latencia de ACK;
\item consumer lag por partición, no solo promedio del grupo;
\item duración y frecuencia de rebalances;
\item tiempo de procesamiento y error rate por operador;
\item tamaño y tiempo de restauración de state stores;
\item registros en DLQ (cola de mensajes no procesables) y edad del record más antiguo;
\item skew de claves y utilización relativa de particiones.
\end{itemize}

Lag es deuda de trabajo, no diagnóstico causal. Puede crecer por falta de CPU, dependencia lenta, hot partition, rebalance, poison pill, GC o throttling.

\subsection{Cómo explicar los conceptos en tres profundidades}

\subsubsection{Log, partición y offset}
\begin{description}
\item[30 segundos.] ``El log es una película retenida; una partición es un carril ordenado; el offset es el señalador de posición dentro de ese carril.''
\item[2 minutos.] Dibuje dos particiones, asigne eventos por key y muestre dos offsets independientes.
\item[5 minutos.] Formalice $P_j=\langle e_{j,0},\ldots\rangle$, explique hashing, skew, retención y replay.
\end{description}

\subsubsection{At-most/at-least/exactly-once}
\begin{description}
\item[30 segundos.] ``El lugar donde confirmo progreso respecto del efecto decide si arriesgo pérdida o repetición.''
\item[2 minutos.] Complete el timeline efecto--offset con un crash en el medio.
\item[5 minutos.] Separe broker, estado local, output Kafka y efecto externo; declare la frontera transaccional.
\end{description}

\subsubsection{Outbox}
\begin{description}
\item[30 segundos.] ``Guardo el cambio y la carta pendiente en la misma transacción; otro proceso envía la carta después.''
\item[2 minutos.] Dibuje tabla de pedidos y outbox; muestre que el relay puede repetir.
\item[5 minutos.] Agregue event ID, ordering por agregado, retries, inbox y limpieza.
\end{description}

\subsubsection{Watermark}
\begin{description}
\item[30 segundos.] ``Es una apuesta informada sobre hasta qué tiempo de evento podemos cerrar resultados.''
\item[2 minutos.] Compare un evento ocurrido a 10:00 y recibido a 10:07 con una ventana ya cerrada.
\item[5 minutos.] Discuta allowed lateness, correcciones, estado y SLO de frescura.
\end{description}

\subsection{Banco de preguntas previsibles y respuestas}

\begin{enumerate}
\item \textbf{¿Kafka es una cola?} Puede usarse para distribución de trabajo, pero conserva un log retenido y permite múltiples grupos/replay; la semántica es más amplia que una cola destructiva clásica.
\item \textbf{¿El offset identifica el evento?} Solo identifica una posición dentro de una partición. Use event ID para identidad de negocio.
\item \textbf{¿Puede haber duplicados aunque el producer sea idempotente?} Sí: comandos de negocio repetidos, reenvíos por aplicaciones, consumidores at-least-once o relays de outbox.
\item \textbf{¿Más particiones siempre mejora throughput?} No: aumenta metadata, conexiones, costo de rebalance, archivos y overhead; además puede perderse la unidad de orden deseada.
\item \textbf{¿Cómo elegir la key?} Primero identifique la unidad que requiere orden y luego valide distribución, skew y capacidad de escalado.
\item \textbf{¿Qué ocurre si cambia el número de particiones?} El mapeo hash puede cambiar y una misma key podría no mantener continuidad histórica según la estrategia; debe planificarse.
\item \textbf{¿Un lag cero prueba exactamente una vez?} No. Solo indica que el marcador alcanzó el final observado; el efecto pudo fallar, duplicarse o ser incorrecto.
\item \textbf{¿Una DLQ resuelve por sí sola los eventos no procesables?} Aísla disponibilidad, pero crea una deuda de corrección y exige ownership, replay seguro y SLO.
\item \textbf{¿Qué limita el paralelismo de un consumer group?} La cantidad de particiones y el costo/estado por partición.
\item \textbf{¿Outbox garantiza orden global?} No. Puede preservar orden por agregado con diseño adicional; múltiples relays y particiones requieren claves y secuencias.
\item \textbf{¿Por qué no usar solo RPC?} RPC puede ser correcto para una respuesta inmediata, pero acopla disponibilidad y tiempo; el log compra desacoplamiento, replay y fan-out a costo de eventualidad y operación.
\item \textbf{¿Watermark igual a reloj global?} No. Es una estimación operativa sobre progreso en event time.
\end{enumerate}


\subsubsection*{Preguntas técnicas adicionales que pueden surgir}

\begin{enumerate}
\item \textbf{¿Por qué el orden se limita a una partición?} Porque cada partición posee una secuencia propia. Coordinar un orden total entre todas las particiones agregaría una autoridad y un cuello de botella global que el diseño evita deliberadamente.
\item \textbf{¿Cambiar el número de particiones puede alterar el mapeo de keys?} Sí. Según el particionador, una misma key puede mapear a otro número cuando cambia $N$. La continuidad histórica y la migración deben planificarse.
\item \textbf{¿Qué diferencia hay entre retención y compactación de log?} Retención elimina segmentos por tiempo o tamaño; log compaction conserva, bajo reglas específicas, el valor más reciente por key. Ninguna equivale a la compactación de SSTables de una LSM estudiada en clase 6.
\item \textbf{¿Qué significa \texttt{acks=all}?} El líder espera las réplicas sincronizadas requeridas por la configuración. La tolerancia real depende también del mínimo de ISR y de las reglas de elección.
\item \textbf{¿Puede perderse un mensaje ya reconocido?} Depende de la política de ACK, replicación, persistencia y failover. Una respuesta del líder solo puede preceder a la copia en seguidores.
\item \textbf{¿Qué pasa con un consumidor que tarda demasiado?} Puede perder su membresía, provocar rebalance y releer desde el committed offset. También puede quedar fuera de la ventana de retención si el atraso es extremo.
\item \textbf{¿El commit de offset debe realizarse por cada registro?} No necesariamente. Confirmar en lotes reduce overhead, pero amplía el rango que puede repetirse tras un crash.
\item \textbf{¿Qué es un poison pill?} Un evento que falla de forma permanente para un consumidor, por esquema inválido, dato imposible o bug. Reintentar sin límite puede bloquear una partición.
\item \textbf{¿CDC y outbox son lo mismo?} No. Outbox es el patrón de persistir explícitamente un evento pendiente con el cambio de negocio; CDC es una técnica para observar cambios del log transaccional y publicarlos.
\item \textbf{¿Qué ocurre si la outbox crece?} Debe medirse edad, cantidad pendiente, reintentos y retención. Una outbox acumulada es evidencia de que el relay o el broker no acompaña la tasa de negocio.
\item \textbf{¿Un watermark descarta automáticamente eventos?} No. El watermark ayuda a disparar cálculos; la política del operador decide descartar, corregir, reabrir o derivar late data.
\item \textbf{¿Cómo se preserva el orden de eventos de una saga?} Use una key por aggregate, secuencia monotónica y consumidores que detecten huecos o versiones incompatibles. El orden entre agregados sigue sin ser global.
\end{enumerate}

\subsection{Ruta esencial y profundización para 180 minutos}

\begin{description}
\item[Núcleo obligatorio.] Orden por partición, key, offset, consumer group, at-most/at-least, frontera exactly-once, outbox/inbox.
\item[Segundo núcleo.] Lag y rebalance como fenómenos operativos; stateful recovery; event time y watermark.
\item[Profundización opcional.] Detalles de ISR, tuning fino, DLQ avanzada y evolución de schemas.
\end{description}

En auditoría, priorice que los estudiantes construyan un timeline correcto y puedan declarar la frontera de la garantía. Cubrir todas las opciones de Kafka sin esa comprensión sería evidencia débil.

\subsection{Checklist personal de estudio}

\begin{itemize}
\item ¿Puedo explicar por qué no existe orden global entre particiones?
\item ¿Puedo distinguir offset, event ID, position y efecto de negocio?
\item ¿Puedo derivar pérdida y duplicación moviendo el crash?
\item ¿Puedo decir con precisión qué cubre el producer idempotente?
\item ¿Puedo nombrar la frontera exacta de una transacción Kafka?
\item ¿Puedo explicar qué deja sin resolver outbox?
\item ¿Puedo recuperar un state store junto con su offset?
\item ¿Puedo explicar un watermark sin llamarlo garantía absoluta?
\item ¿Puedo diagnosticar cuatro causas distintas de lag?
\item ¿Tengo una demo en papel si no hay cluster?
\end{itemize}

\clearpage
\section{Parte C: auditoría, evaluación y evidencias}
\subsection{Preguntas que puede formular un auditor pedagógico}
\subsubsection{¿Por qué comenzar con intuición y no con la definición formal?}

Respuesta sugerida: la intuición activa conocimientos previos y crea una necesidad cognitiva. El caso no reemplaza la definición; la prepara. Inmediatamente después se explicitan modelo, supuestos, garantía y límites de la analogía. Esta secuencia reduce la cantidad de elementos nuevos que el estudiante debe coordinar al mismo tiempo y permite que la formalización responda a una pregunta que ya comprendió.

\subsubsection{¿Cómo verifica que los estudiantes aprendieron y no solo siguieron la narración?}

Respuesta sugerida: cada bloque contiene una predicción individual, una representación que el estudiante debe completar, una decisión con trade-off y una evidencia breve. El diagnóstico inicial reaparece al final, y el exit ticket exige explicar el mecanismo con palabras propias. Las guías alinean los mismos resultados mediante timelines, cálculos y diseños; por tanto, la clase produce artefactos observables, no solo exposición.

\subsubsection{¿Por qué usa casos reales?}

Respuesta sugerida: no se usan como argumento de autoridad ni como entrenamiento de producto. Se usan para mostrar que el mecanismo responde a una restricción real y que tiene métricas, fallas y costos observables. Después del caso se abstrae un método reusable para que el alumno pueda transferirlo a otro dominio.

\subsubsection{¿Cómo evita que una clase de tres horas produzca sobrecarga cognitiva?}

Respuesta sugerida: se divide en ciclos de 12--18 minutos con cambio de actividad; se mantiene un caso conductor y un vocabulario estable; se separa núcleo obligatorio de profundización; se introducen pausas de recuperación y se vuelve varias veces al mapa de la clase. No se intenta ``pasar todas las slides'': se protege la comprensión del núcleo y se deja el detalle avanzado como material de estudio.

\subsubsection{¿Cómo se alinea la clase con la evaluación?}

Respuesta sugerida: cada resultado de aprendizaje aparece en cuatro lugares: una actividad durante la clase, un ejercicio de la guía, una solución con criterio de corrección y una evidencia en el Trabajo Final. La evaluación privilegia la cadena problema--falla--garantía--mecanismo--trade-off--evidencia, coherente con el marco $D=(P,F,G,S,T)$.

\subsubsection{¿Qué hace cuando aparece una concepción errónea?}

Respuesta sugerida: no la corrige únicamente con una afirmación. Solicita una predicción, mueve el punto de falla en un timeline o construye un contraejemplo. Luego pide al estudiante reformular la regla. Las guías docentes identifican errores previsibles y una intervención concreta para cada uno.

\subsubsection{¿Cómo favorece participación sin perder rigor?}

Respuesta sugerida: primero solicita respuesta individual breve para evitar que participen siempre los mismos; después comparación en parejas y finalmente puesta en común. Las preguntas no buscan opinión libre, sino una decisión justificada por garantía, falla y evidencia. Se admiten soluciones distintas cuando preservan los invariantes y explicitan costos.

\subsubsection{¿Qué evidencia conservaría para una auditoría posterior?}

Conservar:
\begin{itemize}
\item resultados anónimos del diagnóstico inicial y final;
\item una muestra de timelines o matrices completadas;
\item respuestas del exit ticket y la devolución realizada en la clase siguiente;
\item una rúbrica aplicada a una actividad;
\item mejoras del Trabajo Final vinculadas con los conceptos de la clase;
\item notas de ajuste docente: qué bloque tomó más tiempo y qué se modificará.
\end{itemize}

\subsubsection{Pregunta específica: ¿por qué enseñar exactly-once si es tan ambiguo?}

Respuesta sugerida: precisamente porque es ambiguo y aparece con frecuencia en industria. La clase obliga a declarar la frontera: entrega, estado Kafka, estado local o efecto externo. El objetivo no es memorizar una etiqueta, sino poder refutar una promesa incompleta.

\subsubsection{Pregunta específica: ¿qué actividad demuestra transferencia?}

Respuesta sugerida: los estudiantes deben diseñar semánticas diferentes para inventario, email y analytics, y luego trasladar el método a telemetría industrial. Si pueden variar key, offset, frontera y SLO según el dominio, la comprensión no quedó atada a CampusShop.

\subsection{Matriz de corrección y preguntas de defensa de la guía}
Use esta sección para convertir las soluciones en una defensa técnica. La respuesta completa debe separar broker, consumidor, estado de procesamiento y efecto de negocio.

\begin{enumerate}
\item \textbf{Log/cola/base.} Debe comparar retención, replay y ownership del progreso. Pregunta oral: ¿quién decide cuándo un record deja de ser necesario?
\item \textbf{Key.} Debe justificar unidad de orden y distribución. Pregunta oral: ¿qué hot key rompería su diseño?
\item \textbf{Particiones.} Debe calcular capacidad con margen y discutir skew/rebalance. Pregunta oral: ¿por qué 16 consumidores no garantizan 16 veces throughput?
\item \textbf{Consumer group.} Debe mostrar asignación antes/después. Pregunta oral: ¿qué consumidores quedan ociosos y por qué?
\item \textbf{Position/commit.} Debe ubicar crash y reinicio. Pregunta oral: ¿qué se relee aunque ya se ejecutó localmente?
\item \textbf{At-most-once.} Debe mostrar la ventana de pérdida y justificar dominio tolerante. Pregunta oral: ¿qué señal detectaría la pérdida?
\item \textbf{At-least-once.} Debe mostrar repetición y deduplicación durable. Pregunta oral: ¿qué sucede si el consumidor cae después del efecto?
\item \textbf{Lag.} Debe medir por partición y edad, no solo promedio. Pregunta oral: ¿cómo distingue hot partition de dependencia lenta?
\item \textbf{Rebalance.} Debe explicar pausa, revocación, restore y cache fría. Pregunta oral: ¿cuándo autoscaling empeora el sistema?
\item \textbf{ACK producer.} Debe declarar falla cubierta y costo de latencia. Pregunta oral: ¿qué sucede si el líder falla antes de replicar?
\item \textbf{Idempotencia producer.} Debe distinguir duplicado técnico y duplicado de negocio. Pregunta oral: ¿qué ID usaría para este último?
\item \textbf{Exactly-once.} Debe nombrar frontera y estado incluido. Pregunta oral: ¿está el email dentro o fuera de esa frontera?
\item \textbf{Kafka Streams.} Debe coordinar offsets, estado y outputs, y explicar abort/replay. Pregunta oral: ¿qué restaura primero?
\item \textbf{Dual write.} Debe construir ambas ventanas de falla. Pregunta oral: ¿qué estado inconsistente produce cada orden?
\item \textbf{Outbox.} Debe incluir relay repetible, event ID, inbox y retención. Pregunta oral: ¿por qué todavía puede publicar dos veces?
\item \textbf{Event time.} Debe asignar el mismo evento a ventanas distintas según reloj. Pregunta oral: ¿qué tiempo representa el negocio?
\item \textbf{Watermark.} Debe justificar tolerancia por costo de error. Pregunta oral: ¿qué hace con datos posteriores al cierre?
\item \textbf{Stateful join.} Debe especificar key, estado, timeout y recuperación. Pregunta oral: ¿qué ocurre si el pago llega luego de expirar el pedido?
\item \textbf{Campaña.} Debe dar semántica diferente a inventario, email y analytics. Pregunta oral: ¿qué subsistema tolera duplicación y cuál no?
\item \textbf{Integración.} Debe incluir experimento reproducible con crash/rebalance. Pregunta oral: ¿qué criterio falsaría su garantía?
\end{enumerate}

\subsection{Errores que deben penalizarse explícitamente}

\begin{itemize}
\item afirmar orden global de un topic particionado;
\item usar offset como identidad de negocio;
\item presentar at-least-once como garantía de efecto exitoso;
\item afirmar exactly-once sin declarar frontera;
\item proponer outbox sin consumidor idempotente;
\item usar lag promedio como único diagnóstico;
\item cerrar ventanas sin política para eventos tardíos.
\end{itemize}

\subsection{Criterio docente de corrección}
La respuesta excelente distingue orden por partición, progreso del consumidor, estado interno y efecto externo. La respuesta insuficiente enumera componentes sin modelar fallas.

\subsection{Evidencias que conviene conservar}

\begin{itemize}
\item una fotografía o exportación del timeline, matriz o diagrama construido por el curso;
\item resultados anónimos del diagnóstico inicial y del exit ticket;
\item una respuesta estudiantil con feedback que muestre revisión conceptual;
\item la conexión entre la actividad del aula y una decisión del Trabajo Final;
\item un registro breve del tiempo real utilizado por bloque y del ajuste realizado para la clase siguiente.
\end{itemize}

\subsection{Checklist del día de la auditoría}

\begin{enumerate}
\item escribir la pregunta central y el mapa de la clase antes de comenzar;
\item pedir una predicción individual antes de abrir la discusión;
\item hacer visible una representación producida por estudiantes;
\item formalizar después de la intuición y explicitar el límite de la analogía;
\item cerrar cada núcleo con problema, garantía, costo y evidencia;
\item proteger al menos ocho minutos para síntesis y exit ticket;
\item recuperar en la clase siguiente una duda real recogida al cierre.
\end{enumerate}

\end{document}
